\section{Background \& Literature Review}

Re-LAION-5B comprises approximately 5.85 billion image-text pairs harvested via Common
Crawl~\cite{Schuhmann_2022_Laion5b}, and was subsequently re-released with improved safety
filtering following the discovery of harmful content in the original LAION-5B
corpus~\cite{Laion_2024_Relaion5B}. The LAION-Aesthetics v2 5+ subset consists of
images filtered for quality using an aesthetic score predicted by a linear classifier
applied to CLIP ViT-L/14 embeddings, and constitutes the training corpus for Stable
Diffusion v1.5~\cite{Schuhmann_2022_LaionAesthetics}. The web-scraped construction
introduces systematic demographic imbalances. Additionally, the absence of pre-computed
embeddings in the re-release requires independently deriving CLIP embeddings for
retrieval. SDv1.5, based on the Latent Diffusion Model framework~\cite{Rombach_2022_CVPR},
performs denoising in a compressed latent space via a DDPM U-Net conditioned on CLIP
text embeddings~\cite{Ho_2020_DDPM, Radford_2021_CLIP}. Higher classifier-free guidance
scales amplify demographic stereotypes by pushing generation more strongly toward the
model's learned conditional distribution~\cite{Ho_2022_CFG}. SDv1.5's fully documented
training lineage makes it uniquely suited for mechanistic bias auditing; later models
such as SDXL~\cite{Podell_2023_SDXL} rely on proprietary corpora that preclude
independent verification.

Spurious features are low-level, non-semantic image characteristics that correlate with
sensitive attributes. They manifest across the full visual hierarchy, including pixel
colour statistics, frequency signatures, background scene context, object co-occurrence
patterns, and structural geometry. Zeng et al.~\cite{Zeng_2024_UnderstandingBiasinLargeScaleVisualDatasets}
operationalise this through a transformation-based classification framework in which
each transformation, spanning pixel shuffle, patch shuffle, mean RGB, frequency
filtering, Canny edge, SAM contour, monocular depth, semantic segmentation, object
detection, and image captioning, isolates a distinct visual property while discarding
all others. A dataset classifier trained on each transformed representation quantifies
how much discriminative signal that property carries independently, revealing that
dataset identity is redundantly encoded across colour, texture, structural geometry,
and semantic content simultaneously. Crucially, Zeng et al.\ further demonstrate that
generative models trained on biased datasets reproduce those biases in their outputs,
directly motivating the present study of Re-LAION-5B and SDv1.5. Meister et
al.~\cite{Meister_2023_GenderArtifactsinVisualDatasets} extend this perspective to
demographic attributes by introducing a systematic occlusion-based probing protocol
that progressively removes the person, background, and contextual objects from images
to isolate which regions carry gender-correlated signal. By training a classifier
separately on the full image, person-only, background-only, silhouette, pose keypoints,
and masked variants, they establish that gender artifacts are distributed across the
entire image structure rather than localised to the subject, undermining
fairness-through-blindness assumptions. This thesis adapts both frameworks. Zeng et
al.'s transformation hierarchy is extended to a three-class Re-LAION-5B / SDv1.5 /
COCO setting, and Meister et al.'s occlusion protocol is applied to probe demographic
artifact distributions in both the retrieved and generated corpora.

Occupation prompts, drawn systematically from the ISCO-08 taxonomy~\cite{ILO_ISCO08},
provide an interpretable probe for demographic priors, as neutral prompts reliably
surface representational biases~\cite{Naik_2023_SocialBiasesThroughTheTextToImageGenerationLens,
	Bianchi_2023_TTIGenAmplifiesDemographicStereotypes,
	Luccioni_2023_AnalyzingSocietalRepresentationsDiffusionModels}. Demographic annotation
at web scale relies on automated classifiers applied to segmented facial skin regions
to prevent background shortcut learning~\cite{Matias_2024_EnhancingFairnessMST}, with
skin tone operationalised via the ten-point Monk Skin Tone scale~\cite{Schumann_2024_MonkSkinTone}
and evaluated against CCv2~\cite{Porgali_2023_CCv2} and FACET~\cite{Gustafson_2023_FACET}
benchmarks. Automated annotation introduces structured measurement error that can
produce artificial spurious correlations~\cite{Buolamwini_2018_GenderShades}.
Inference-time debiasing methods operate on frozen models without retraining. This
thesis employs ITI-GEN~\cite{Zhang_2023_ITIGEN}, which learns inclusive text-space
tokens grounded in demographic reference images, composed with ControlNet
OpenPose~\cite{LZhang_2023_ControlNet} to constrain body pose structure, and
Chamfer-based colour alignment~\cite{Shum_2025_ColorAlignment} to neutralise
colour-level spurious priors across the denoising trajectory.